% Cote 119: offprints, reports, courses, CV and letters (typed and printed part): archive cote 119, site page 295 (PDF page 296).
% typescript. Draft compared: none.
% AI-transcribed by Claude Opus 5.5 (claude-opus-5-5, adaptive thinking, effort high), Message Batches API, batch r119_full (msgbatch_01DwZNm5Y7YsKBWyAN4sFs8w), 2026-11-13; not yet checked by eye.
\scanpage{119/295}
\mspage{- 8 -}

de $H^1(k, G)$. Supposons que $x$ et $y$ aient même images dans $H^1(k', G)$
et dans $H^1(k'', G)$ o\over{u}{ù} $k'$ et $k''$ sont deux extensions finies de $k$ de
degrés premiers entre eux (par exemple $\hand{[}k':k\hand{]} = 2$ et $\hand{[}k'':k\hand{]} = 3$). Ceci
\uline{n'entraîne pas} $x = y$ contrairement à ce qui se passe dans le cas abélien;
on peut en construire des exemples, en prenant $G$ non connexe; j'ignore
ce qu'il en est lorsque $G$ est connexe.

§5. Les formes traces

Il s'agit de la structure de la forme quadratique $Tr(x^2)$ associée à
une $k$-algèbre de dimension finie. Deux cas particuliers ont été considérés :

5.1. Algèbres centrales simples

Soit $A$ une telle algèbre, supposée de degré fini $n^2$ sur $k$. On lui
associe la forme quadratique $q_A$ définie par
\[ q_A(x) = Trd_{A/k}(x^2). \]
Notons $q^\circ_A$ la forme trace associée à l'algèbre de matrices $M_n(k)$ de
même rang que $A$; c'est la somme directe d'une forme hyperbolique de rang
$n(n-1)$ et d'une forme unité $\hand{\langle}1,1,\ldots,1\hand{\rangle}$ de rang $n$.

On désire comparer $q_A$ et $q^\circ_A$ . Il y a deux cas à distinguer :

(5.1.1) $n$ \uline{est impair}.

Les formes $q_A$ et $q^\circ_A$ sont alors isomorphes; cela résulte du théorème de
Springer cité en (4.3).

(5.1.2) $n$ \uline{est pair}.

Soit $(A)$ la classe de $A$ dans le groupe de Brauer de $k$. Le produit
de $(A)$ par l'entier $n/2$ est un élément $a$ de $Br_2(k) = H^2(k, Z/2Z)$. On a :
\[ w_1(q_A) = w_1(q^\circ_A) \qquad \text{et} \qquad w_2(q_A) = w_2(q^\circ_A) + a . \]
(La formule relative à $w_1$ est facile. Celle relative à $w_2$ s'obtient en
considérant l'homomorphisme $PGL_n \mathrel{-\!\!>} SO_{n^2}$ donné par la représentation
adjointe et en montrant, par un calcul de poids et racines, que cet
homomorphisme ne se relève pas au groupe $Spin_{n^2}$ si $n$ est pair.)

5.2. Algèbres commutatives étales

Soit $E$ une telle algèbre, soit $n$ son rang et soit $q_E$ la forme
trace correspondante. Les invariants $w_1$ et $w_2$ de $q_E$ sont donnés par
une formule connue (\uline{Comm}.\uline{Math}.\uline{Helv}. 59, 1984). Le cours a donné une
démonstration de cette formule quelque peu différente de la démonstration
originale, et a appliqué le résultat obtenu aux équations quintiques à la
Kronecker-Hermite-Klein.

\mspage{\hand{295}}
